\begin{theorem}[Unramified norm, trace, and conductor compatibility]\label{mod:ramification-unramifiedcompatibility}
Let $K/F$ be a finite valuative extension of nonarchimedean local fields and
assume that its ramification index is one.  Put
\[
 f=[K:F]=[k_K:k_F].
\]
The integral trace and norm are defined as maps
\(
 \operatorname{Tr}_{\mathcal O}\colon\mathcal O_K\to\mathcal O_F
\)
and
\(
 N_{\mathcal O}\colon\mathcal O_K\to\mathcal O_F
\), and their coercions to the fraction fields are the field trace and norm.

Every uniformizer $\pi$ of $F$ remains a uniformizer after mapping to $K$, and
\[
 N_{K/F}(\pi)=\pi^f.
\]
The residue extension has degree $f$.  For every $x\in\mathcal O_K$, reduction
commutes with both integral maps:
\[
 \overline{\operatorname{Tr}_{\mathcal O}(x)}
   =\operatorname{Tr}_{k_K/k_F}(\bar x),
 \qquad
 \overline{N_{\mathcal O}(x)}=N_{k_K/k_F}(\bar x).
\]
The norm assertion is proved for every integral element, and hence in
particular for every integral unit required by the manuscript.

For every $a\in\mathbf Z$, the trace image of the full fractional lattice is
exactly the lattice of the same depth.  Lean states this without set-image
abbreviations as
\[
 y\in\mathfrak p_F^a
 \quad\Longleftrightarrow\quad
 \exists x\in\mathfrak p_K^a,\ \operatorname{Tr}_{K/F}(x)=y.
\]
For every $m\in\mathbf N$, including $m=0$, the corresponding norm statement is
\[
 u\in U_F^m
 \quad\Longleftrightarrow\quad
 \exists x\in U_K^m,\ N_{K/F}(x)=u.
\]
Here the normalization is exactly $U^0=\mathcal O^\times$ and
$U^m=1+\mathfrak p^m$ for $m>0$.  Positive-depth surjectivity is obtained by
the graded trace correction and completeness; depth zero is obtained from
finite-residue-field norm surjectivity followed by the $U^1$ result.

Consequently triviality is preserved at every individual lattice or unit
layer.  In particular, with the manuscript's sign convention that
$\mathfrak p^{-n}$ is the largest additive triviality lattice,
\[
 \operatorname{IsAdditiveConductor}_K
   (\psi\circ\operatorname{Tr}_{K/F},n)
 \quad\Longleftrightarrow\quad
 \operatorname{IsAdditiveConductor}_F(\psi,n),
\]
and
\[
 \operatorname{IsMultiplicativeConductor}_K
   (\chi\circ N_{K/F},m)
 \quad\Longleftrightarrow\quad
 \operatorname{IsMultiplicativeConductor}_F(\chi,m).
\]

Finally, for a uniformizer $\pi\in F^\times$, Lean proves the elementwise norm
image formula
\[
 u\in N_{K/F}(K^\times)
 \quad\Longleftrightarrow\quad
 \exists a\in\mathbf Z\ \exists\varepsilon\in U_F^0,
 \quad u=\pi^{fa}\varepsilon.
\]
It also constructs the explicit multiplicative equivalence
\[
 F^\times/N_{K/F}(K^\times)
   \simeq \operatorname{Multiplicative}(\mathbf Z/f\mathbf Z).
\]
Thus the quotient has cardinality $f$, is cyclic, the class of $\pi$ maps to
$1\bmod f$, and the subgroup generated by that class is the whole quotient.
\LeanFile{LanglandsFirstMainLemma/Ramification/UnramifiedCompatibility.lean}
\lean{LanglandsFirstMainLemma.integralTrace}
\lean{LanglandsFirstMainLemma.integralNorm}
\lean{LanglandsFirstMainLemma.unramified_uniformizer_norm}
\lean{LanglandsFirstMainLemma.unramified_residue_degree}
\lean{LanglandsFirstMainLemma.unramified_residue_trace_norm}
\lean{LanglandsFirstMainLemma.unramified_trace_lattice_image}
\lean{LanglandsFirstMainLemma.unramified_norm_unitFiltration_image}
\lean{LanglandsFirstMainLemma.unramified_addCharTrivialOnLattice_iff}
\lean{LanglandsFirstMainLemma.unramified_quasiCharTrivialOnUnitFiltration_iff}
\lean{LanglandsFirstMainLemma.unramified_additiveConductor_compTrace}
\lean{LanglandsFirstMainLemma.unramified_multiplicativeConductor_compNorm}
\lean{LanglandsFirstMainLemma.mem_unramified_norm_range_iff_uniformizer}
\lean{LanglandsFirstMainLemma.unramifiedNormQuotientEquiv}
\lean{LanglandsFirstMainLemma.unramifiedNormQuotient_natCard}
\lean{LanglandsFirstMainLemma.unramifiedNormQuotient_isCyclic}
\lean{LanglandsFirstMainLemma.unramifiedNormQuotientEquiv_uniformizer}
\lean{LanglandsFirstMainLemma.unramifiedNormQuotient_uniformizer_generates}
\leanok
\SourceText{Lemma lem:unramified-compatibility and Proposition prop:unramified-new.}
\uses{mod:localfield-extension,mod:localfield-residuefield,mod:basic-characterconductors}
\FormalizationContract
\end{theorem}
